%!TEX TS-program = xelatex
%!TEX encoding = UTF-8 Unicode

\documentclass[a5paper,twosided,11pt]{book}
\usepackage[papersize={140mm,210mm},hmargin=9mm,tmargin=10mm,bmargin=20mm]{geometry}




\usepackage{calc}
\usepackage{amsmath,xltxtra,xunicode}
\usepackage{titlesec}
\usepackage{fontspec}
\usepackage{wasysym,oz,pxfonts,txfonts}


\usepackage{xeCJK}
\XeTeXlinebreaklocale "zh" 
\XeTeXlinebreakskip = 0pt plus 1pt 

% =========================================
\usepackage[none]{hyphenat}
\usepackage{soul}
\usepackage[hidelinks]{hyperref}
\usepackage{lipsum}
\usepackage{lastpage}
\usepackage{everypage}
\usepackage{fancyhdr}
\usepackage{listings}
\usepackage{paralist}
\usepackage{enumerate}
\usepackage{enumitem}
\usepackage{tocloft}
\usepackage{longtable}
\usepackage{tabu}
\usepackage{makecell}
\usepackage[dvipsnames]{xcolor}
\usepackage{svg}
\usepackage{graphicx}
\usepackage{tikz}
\usepackage{eso-pic}
\usepackage{fontawesome5}
\usepackage{multirow}
\usepackage{multicol}
\usepackage{ragged2e}
\usepackage{tcolorbox}
\setmainfont{Source Serif Pro}
\setromanfont{Source Serif Pro}
\setromanfont{Inria Serif}
\setsansfont{IBM Plex Sans}
\setmonofont{JetBrains Mono NL}
\setCJKmainfont{Source Han Serif CN}
\setCJKromanfont{Source Han Serif CN}
\setCJKsansfont{Source Han Sans CN}
\setCJKmonofont{Source Han Sans CN}
\title{AOSC OS/Retro: User Manual}
\author{Various Authors at AOSC}
\date{2021-09-20}
\pagestyle{plain}
% =========================================

% =========================================
% START CORPORATION CONFIG
\setlength{\parindent}{0pt}
\setlength{\parskip}{4pt}
\setlength{\baselineskip}{15pt}
\tabulinesep=6pt
% END CORPORATION CONFIG
% =========================================

% =========================================
% START COMMON STYLESHEET
% Table
\renewcommand\theadalign{l}
\renewcommand\theadfont{\sffamily\bfseries}
% Font Size

% END COMMON STYLESHEET
% =========================================


\graphicspath{{/home/neruthes/EWS/AOSC/optical-media-box/_img}}


% Headers & Footers
\fancypagestyle{plain}{
    \renewcommand{\headrulewidth}{0pt}
    \renewcommand{\footrulewidth}{0.4pt}
    \fancyhf{}
    \fancyfoot{}
    \fancyfoot[LE,RO]{\small\ttfamily\thepage}
    \fancyfoot[LO]{\small\sffamily{AOSC OS/Retro: User Manual}}
    \fancyfoot[RE]{\small\sffamily{AOSC OS/Retro 用户手册}}
}


% CoverPage
\newcommand{\docversion}[0]{0.1.1-pre1}

\newcommand{\makecoverpage}[0]{{

    % \AddToShipoutPictureBG*{%
    %     % Two major boxes
    %     \put(25mm,50mm){%
    %         \includegraphics[width=90mm]{retrologo-skeleton.pdf}
    %     }%        
    % }

    % Start real content
    \raggedbottom
    \RaggedRight
    \pagestyle{empty}
    \fontspec{Source Han Serif CN Heavy}
    \CJKfontspec{Source Han Serif CN Heavy}

    \center

    \color{white}
    \pagecolor{black}


    % \hspace{1pt}\vspace{1pt}

    \includegraphics[width=24mm,height=24mm]{retro-logo.png}
    \vspace{4pt}

    {
        % \bfseries
        \huge
        \sffamily
        \fontspec{Noto Sans SemiCondensed SemiBold}
        \CJKfontspec{Noto Sans CJK SC Medium}
        
        USER MANUAL

        用户手册
    }

    

    % \vspace{10pt}{
    %     \rmfamily\color{white!82!black}\normalsize\mdseries

    %     Version \docversion\\
    %     September 2021
    % }

    \vfill
    \includegraphics[width=70mm]{anan-mac.png}
    \vfill

    {

        % \begin{minipage}{\textwidth}
            % \center
            % {\hfill
            % \includegraphics[width=14mm,height=14mm]{aosc.pdf}\hspace{10pt}
            % \includegraphics[width=24mm,height=24mm]{retro-logo.png}
            % \hfill}
        % \end{minipage}
        % \vspace{5pt}
        
        \fontspec{Source Han Serif CN Heavy}
        \bfseries
        \LARGE
        AOSC OS/Retro

        \vspace{5pt}
        {
            \rmfamily
            \normalsize\mdseries
            Enjoying AOSC OS on Your Retro Devices
        }
    }
    
    

    \pagebreak
    \color{black}
    \pagecolor{white}
    \pagebreak

    % \cleardoublepage
    \setcounter{page}{1}
    \rmfamily\mdseries\normalsize

    
}}

\newcommand{\makebackcoverpage}[0]{{
    \hspace{1pt}\vfill

    \small\sffamily\mdseries
    \baselineskip=16pt
        {\tpVersion}{ }{\docversion}\\
        \tpCopyright\\
        \tpLicense{\underline{CC BY 4.0}}
}}







\setulcolor{gray}

\newcommand{\icode}[1]{{\ttfamily{\ul{#1}}}}

\newcommand{\codeblock}[1]{
\begin{tcolorbox}[colframe=black,colback=white]
\ttfamily\small
#1
\end{tcolorbox}
}

\newcommand{\codenote}[1]{
    {\footnotesize #1}
}

\renewcommand{\chapter}[1]{
    \pagebreak
    \refstepcounter{chapter}
    \addcontentsline{toc}{chapter}{\protect\numberline{\thechapter}#1}
    \chaptermark{#1}
    {
        \sffamily
        % \rmfamily
        % \fontspec{IBM Plex Sans}
        \huge\bfseries
        \scshape
        {\thechapter}. {#1}
    }\par\vspace{20pt}
}

\begin{document}
    \baselineskip=16pt
    \sloppy
    \raggedbottom
    \RaggedRight
    \overfullrule=10pt

    % \makecoverpage

    \setcounter{tocdepth}{1}
    \sffamily{\tableofcontents}
    % \clearpage

    % \hspace{40mm}

    % \large
    % \vspace{40mm}
    % Now that we have introduced our rationale to maintaining a Retro branch for AOSC OS
    % (we will now refer to this as AOSC OS/Retro), this page will serve as an introduction
    % to our design specifications and goals in maintaining this branch.
    % \vfill

    \pagestyle{plain}
    \rmfamily\mdseries\normalsize

    \chapter{Specifications}

    Now that we have introduced our rationale to maintaining a Retro branch for AOSC OS
    (we will now refer to this as AOSC OS/Retro), this page will serve as an introduction
    to our design specifications and goals in maintaining this branch.

    In terms of end-user experience and management characteristics, AOSC OS/Retro is a standard AOSC OS distribution.
    However, changes are made in terms of dependencies, features, and maintenance schedules
    In brief, AOSC OS/Retro will\ldots

    \begin{enumerate}
        \item Come with limited architectural support - only those specified
            as Retro will be maintained in accordance to the Retro specifications.
        \item Share the same package tree as the mainline distribution, as well as the Core.
        \item Share the same set of maintenance tools as the mainline distribution.
        \item Strip down package features in interest of conserving storage and memory.
        \item Ship in different flavours as otherwise would with the mainline distributions.
        \item Update on a slower schedule, with a few exceptions.
    \end{enumerate}

    Now, we will present the commonalities and differences in detail.




    \section{Target Hardware}
    
    AOSC OS/Retro currently support the following processing architectures and devices\ldots
    \pagebreak

    \begin{tabu} {XX}
        \hline
        {Architecture} &
        {Typical Devices}\\
        \hline
        {ARMv4, soft float with no Thumb support}
            & {HP Jornada 710/720/728/820, etc.}\\
        {ARMv6, hard float}
            & {Raspberry Pi 1, etc.}\\
        {ARMv7, hard float with NEON SIMD support}
            & {Raspberry Pi 2, etc.}\\
        {Intel 80486 compatible or above (\icode{i486})}
            & {IBM Palm Top PC110, Toshiba Satellite 430CDT, Dell Latitude CPi, Asus EeePC 901 etc.}\\
        {Loongson 2F (\icode{loongson2f})}
            & {Lemote YeeLoong 8089B/D, etc.}\\
        {PowerPC G3 and above, 32-bit big endian (\icode{powerpc})}
            & {Apple iBook G3, Apple iMac G4, etc.}\\
        {PowerPC G5 and above, 64-bit big endain (\icode{ppc64})}
            & {Apple PowerMac G5, IBM eServer pSeries 610, etc.}\\
        \hline
    \end{tabu}

    \section{Common Components and Tools}
    
    As AOSC OS/Retro is \textit{not} a fork but a \textit{branch} of the mainline AOSC OS distribution,
    it will share not only the package tree, but also the Core and the maintenance tools as the mainline\ldots
    
    \begin{itemize}
    \item Package Tree: \textbf{aosc-os-abbs}, on the \icode{retro} branch.
        \begin{itemize}
            \item This means that AOSC OS/Retro will also use systemd as the default init system.
                {\itshape Now before you start yelling, systemd runs fine on an IBM Palm Top PC110
                with a 33MHz 80486SL processor, 20MiB of RAM, and a 2GiB CF card. We've tested it to be sure.}
        \end{itemize}
    \item AOSC OS Core will be shared, but only updates based on what are synced to the retro branch. See Maintenance Schedule.
    \item Packaging/Maintenance Tools\ldots
        \begin{itemize}
            \item Autobuild3, for automatic packaging from Autobuild3 manifests.
            \item ACBS, for tree-based Autobuild3 manifest management and packaging.
            \item Ciel, for containerised packaging.
            \item Various scriptlets, shared with AOSC OS.
        \end{itemize}
    \end{itemize}
    

    \section{Dependencies}

    As AOSC OS/Retro targets storage and performance constraint devices, unlike its mainline counterpart,
    AOSC OS/Retro will ship packages with minimal optional feature enabled.
    Listed below are a few general rules to be followed\ldots

    \begin{itemize}
        \item Base distribution (containing \icode{admin-base}, \icode{boot-base}, \icode{core-base}, \icode{editor-base},
            \icode{kernel-base}, \icode{network-base}, \icode{systemd-base}, \icode{util-base}, and \icode{web-base})
            must not introduce dependency to Python (\icode{python-2}, \icode{python-3}).
        \item Language bindings (Java, Perl, Python, etc.) must not be enabled by default,
            unless required by another package essential to the distribution.
        \item Glibc is \textit{only} to ship with the \icode{C} and \icode{C.UTF-8} locales pre-generated,
            with others generated by user configuration (defined in \icode{/etc/locale.gen} but commented by default).
        \item Strip down \textit{all optional} dependencies, unless such package is from the Core,
            or otherwise discussed on a case-by-case basis.
        \item \textit{All packages} are to be built with Link-Time Optimisation enabled,
            unless such optimisation leads to build failure (to be reported to upstream).
        \item Non-performance critical applications are to be built with the \icode{-Os}
            (\icode{AB\_FLAGS\_OS=1} in \icode{autobuild/defines}) optimisation level to conserve space.
        \item Manpages and Texinfo documents will be shipped, but all other forms of documentation
            (examples, HTML, gtk-doc, etc.) will be omitted.
        \item Linux Kernel must boot \textit{without} Dracut, unless used on a RAID setup.
            Dracut is \textit{not} shipped with default distributions.
    \end{itemize}

    \section{Distribution Features}

    AOSC OS/Retro will ship in several flavours, Base, X11, Server, and Trinity.
    All flavours will come fully localised (once locale is enabled by the end-user) as well as a generic,
    non-optimised Linux Kernel for each architecture.

    \begin{itemize}
        \item The Base flavour contains a minimal bootable and non-graphical system with tools essential to system management, basic text editing (GNU Nano), networking (wired and wireless), and basic user functions (compression, documentation browser, pager, power conservation, system monitoring, etc.).
        \item The X11 flavour contains a minimal bootable and graphical system, with all components listed above for the Base flavour, added with an X11-based desktop environment and other graphical utilities.
        \item The Server flavour contains a minimal bootable and server-oriented system, with all components listed above for the Base flavour, added with tools and utilities for server-oriented applications.
        \item The Trinity flavour contains a fully-featured bootable and graphical system, with all components listed above for the Base flavour, added with a fully-featured desktop interface powered by the Trinity Desktop Environment.
        \item Network management on all flavours will be provided with NetworkManager, and graphical systems will come with NetSurf as the default Web browser.
    \end{itemize}

    Extra packages, such as Firefox and more feature-complete desktop environments will be available from the community repository, however, hardware requirement checks will be enforced based on processor and memory installed on your AOSC OS/Retro device (i.e., package installation will be aborted when attempting to install Firefox on a computer without SSE2 SIMD support).

    \section{Maintenance Schedules}

    AOSC OS/Retro will be maintained on the \icode{retro} branch,
    sharing the same package tree with the mainline distribution.
    However, in interest of both the maintainer's reasonable maintenance effort,
    as well as the longevity and usability of the target devices,
    AOSC OS/Retro will update on an \textit{annual schedule}.

    After the first update cycle of a year, the retro branch will perform a dual-direction
    merge with the \icode{stable} branch from the mainline distribution (\icode{stable} => \icode{retro}, then \icode{retro} => \icode{stable}).
    After which, no further merge or reverse merge will be allowed.
    Package versions in the \icode{retro} branch will remain constant unless\ldots

    \begin{itemize}
        \item Any bugfixes are found to be necessary by users or developers.
        \item A security update is made available that requires a version update.
            If necessary, changes could be cherry-picked from the \icode{stable} branch.
    \end{itemize}

    At the end of each annual cycle, a new distribution tarball will be made available on the downloads page, as well as an update CD image containing a local repository containing all system updates. A full AOSC OS/Retro repository will also be provided in forms of a tarball or a set of CD/DVD image.





    \chapter{Goals}

    AOSC OS/Retro will be maintained with a few goals in mind, relating to system performance, storage requirements, and peripheral support. This chapter will also serve to outline AOSC OS/Retro's system requirements.

    This chapter will then be split into sections, containing requirements and metrics shared and specific to each of our target architectures.

    \section{Common Metrics}

    \begin{itemize}
        \item AOSC OS/Retro's Base flavour should install onto a 810MiB hard disk drive (commonly found with 486-class systems), while the X11 flavour should install onto a 1.2GB drive (commonly found on Intel x86 computers from ~1996).
            \begin{itemize}
                \item After the system is installed, there should be enough space for a 64MiB swap area and future system updates (assuming one package is cached onto the hard disk at a time, using the update CD).
                \item Users should expect to conserve ~100MiB of hard disk space for network- or internet-based system updates.
            \end{itemize}
        \item AOSC OS/Retro should not require any form of network access for normal usage, assuming the user has obtained a copy of local repository.
        \item AOSC OS/Retro should support common ISA/EISA (or PCMCIA), PCI (or CardBus), PCI Express (or ExpressCard), SCSI, as well as USB, PS/2, Serial and Parallel peripherals.
        \item AOSC OS/Retro should support dial-up, 10/100/1000Mbps Ethernet, as well as 802.11a/b/g/n/ac/ax wireless connections.
        \item AOSC OS/Retro should boot from IDE/EIDE/CE-ATA/SATA/SCSI-based hard disk drives. AOSC OS/Retro may boot from USB, optical media, or other forms of external/removable storage, but this will not officially supported.
    \end{itemize}







    \chapter{System Requirements}

    \section{General Advice}

    In order to ensure an acceptable and productive AOSC OS/Retro experience, we recommend the following when considering installation of AOSC OS/Retro on your devices:

    \begin{itemize}
        \item Minimum requirements are here for technicality's sake,
            it does not necessarily mean that AOSC OS/Retro would run with sufficient performance
            on said configuration. We have included specific advice for each architecture.
        \item Flash-based storage devices, such as bridged CompactFlash cards
            (connected via IDE, SCSI, etc.) will significantly improve system performance on devices from the 1990s.
            CompactFlash cards and adapters are affordable, consumes less power,
            and lowers overall system heat output when compared with the older,
            (now) frequently failing mechanical IDE hard disk drives.
        \item Control your expectations. While AOSC OS/Retro runs well on computers from well over 20 years ago,
            they won't run Crysis or even play YouTube videos. We advise that you plan your applications wisely,
            especially on devices that are severely performance-limited.
    \end{itemize}


    \section{x86/i486 Systems}

    \subsection{Base \& Server}

    On the 32-bit x86 architecture, AOSC OS/Retro Base and Server requires the following system components\ldots

    \begin{compactitem}
        \item Processor: Intel 80486 or compatible, FPU (Floating Point Unit) not required.
        \item System Bus: ISA, EISA, PCI, or PCI Express based system devices. MCA (Micro Channel Architecture) not supported.
        \item RAM: 12MiB (32MiB swap).
        \item Storage: 810MB ({\textasciitilde}514MiB).
            \begin{compactitem}
                \item (Ultra) DMA via PCI Bus Mastering will significantly improve system performance.
            \end{compactitem}
        \item Input Device: PS/2 or Serial Port Keyboard. Mouse not required.
        \item Display: VGA or compatible, or serial terminal.
    \end{compactitem}

    \subsection{X11}

    AOSC OS/Retro X11 requires the following system components\ldots

    \begin{compactitem}
    \item Processor: Intel 80486 or compatible, FPU (Floating Point Unit) not required.
        \begin{compactitem}
            \item Intel Pentium II 233MHz, AMD K6, Cyrix MediaGX, Via C7 or above will significantly improve graphical experience.
            \item Intel Pentium III 500MHz, AMD K6-II/III or above recommended for video playback using MPlayer.
        \end{compactitem}
    \item System Bus: ISA, EISA, PCI, or PCI Express based system devices. MCA (Micro Channel Architecture) not supported.
    \item RAM: 48MiB (64MiB swap).
        \begin{compactitem}
            \item 128MiB or above recommended for Internet browsing.
        \end{compactitem}
    \item Storage: 1.2GB (~1141MiB).
        \begin{compactitem}
            \item 4.0GB (~3814MiB) recommended for local multimedia storage.
            \item (Ultra) DMA via PCI Bus Mastering will significantly improve system performance.
        \end{compactitem}
    \item Input Device: PS/2 or Serial Port Keyboard and Mouse.
        \begin{compactitem}
            \item Touchscreen will be supported via I2C, Serial, or USB connections.
        \end{compactitem}
    \item Display: VGA or compatible.
        \begin{compactitem}
            \item ISA/EISA video cards are not recommended, VESA Local Bus will significantly improve video performance.
            \item PCI and PCI Express video cards recommended, especially those with OpenGL 2.1 support (often found after ~2002), as this will allow for GPU-based video playback acceleration.
        \end{compactitem}
    \end{compactitem}

    \subsection{Trinity}

    AOSC OS/Retro Trinity requires the following system components\ldots

    \begin{compactitem}
    \item Processor: Intel Pentium II 233MHz, AMD K6, Cyrix MediaGX, Via C7 or above.
        \begin{compactitem}
        \item Intel Pentium III 500MHz, AMD K6-II/III or above recommended.
        \end{compactitem}
    \item System Bus: ISA, EISA, PCI, or PCI Express based system devices. MCA (Micro Channel Architecture) not supported.
    \item RAM: 256MiB (256MiB swap).
        \begin{compactitem}
        \item 512MiB or above recommended for Internet browsing.
        \end{compactitem}
    \item Storage: 3.2GB (~3052MiB).
        \begin{compactitem}
        \item 8GB (~7629MiB) recommended for local multimedia storage.
        \item (Ultra) DMA via PCI Bus Mastering will significantly improve system performance.
        \end{compactitem}
    \item Input Device: PS/2 or Serial Port Keyboard and Mouse.
        \begin{compactitem}
        \item Touchscreen will be supported via I2C, Serial, or USB connections.
        \end{compactitem}
    \item Display: VGA or compatible.
        \begin{compactitem}
        \item PCI and PCI Express video cards recommended, especially those with OpenGL 2.1 support (often found after ~2002), as this will allow for GPU-based video playback acceleration.
        \end{compactitem}
    \end{compactitem}


    \section{PowerPC Systems}

    All AOSC OS/Retro flavours should run on all supported devices on this architecture - that is, PowerPC-based Apple Macintosh computers with New World ROM support. If you would like to run Trinity Desktop Environment on your PowerPC devices, you are advised to have at least 256MiB of RAM installed.

    \begin{compactitem}
    \item Portables...
        \begin{compactitem}
        \item PowerBook G3 "Lombard" and "Pismo" models.
        \item All iBook G3, iBook G4, PowerBook G4 models.
        \end{compactitem}
    \item Desktops...
        \begin{compactitem}
        \item Power Macintosh G3 "Blue and White" models.
        \item All Power Macintosh G4 and G5 models.
        \item All iMac G3 and G4 models.
        \item All eMac models.
        \item All G4-based Mac Mini models.
        \item All G4- and G5-based Xserve models.
        \end{compactitem}
    \end{compactitem}

    \section{Loongson 2F Systems}

    All AOSC OS/Retro flavours should run reasonably well on Loongson 2F-based devices
    (with performance close to an Intel Pentium III Coppermine).
    Currently, we only support the Lemote YeeLoong 8089 family of laptops.

    For storage requirements, you may refer to the x86 section above.

    \section{ARMv4 Systems}

    For ARMv4, AOSC OS/Retro supports a limited range of devices.
    Specifically the HP Jornada 710/720/728/820 Windows CE handheld computers.
    AOSC OS/Retro Base and Server should run on any of the supported devices with their standard RAM capacity (32MiB or 64MiB),
    X11 requires that you install a 64MiB RAM capacity ROM card to ensure acceptable performance.
    Trinity Desktop Environment is not supported, for obvious performance reasons.

    For storage requirements, you may refer to the x86 section above.

    \section{ARMv6 Systems}

    For ARMv6, AOSC OS/Retro currently support only the Raspberry Pi 1 models
    with 256MiB - 512MiB of RAM. All AOSC OS/Retro flavours should run reasonably well on these devices.

    For storage requirements, you may refer to the x86 section above.

    \section{ARMv7 Systems}

    For ARMv7, AOSC OS/Retro currently support only the Raspberry Pi 2 models,
    with Allwinner device support planned.
    All AOSC OS/Retro flavours should run reasonably well on these devices.

    For storage requirements, you may refer to the x86 section above.

    % \chapter{Introduction}
    % \section{114514}






    \chapter{Installation (Generic)}

    This guide appies to i486, Loongson 2F, and PowerPC (32/64).

    Apart from guides on how to create LiveKit media, this guide also outlines two main use cases for LiveKit,
    the installation and maintenance of AOSC OS.
    LiveKit is a portable AOSC OS installation (i.e. a Live environment) which provides a self-contained
    and complete environment deployed on media such as USB flash drives and optical discs.


    \section{Creating the LiveKit Media}
    
    As you read this manual, you should have received such discs in the same box.
    If not, you may create one.

    Creating LiveKit media is easy. This section shows how to create USB LiveKit media under Linux
    (or other *nix operating systems with the dd command available), Windows, and macOS.
    To create an optical media, simply use any optical disc burner to write the .iso you have obtained from the Downloads page.

    \subsection{Linux, *nix, and macOS}

    In operating systems using the \icode{/dev/sdXY} (where X is the device number, denoted with a letter;
    Y is the partition number, denoted with a number), please first identify your USB device (via \icode{lsblk} or simply \icode{ls /dev/sd*}),
    then execute the following commands \textit{as root}:

\codeblock{
dd if=aosc-os\_livekit\_DATE\_ARCH.iso \textbackslash\\
    of=/dev/sdXY status=progress\\
sync
}
% \begin{tcolorbox}[colframe=black,colback=white]
% \small\begin{verbatim}
% dd if=aosc-os_livekit_DATE_ARCH.iso \
%     of=/dev/sdXY status=progress
% sync
% \end{verbatim}
% \end{tcolorbox}

    \codenote{NOTE: aosc-os\_livekit\_\$DATE\_\$ARCH.iso, where \$DATE is the release date, and \$ARCH
        is the architecture number which the LiveKit targets. Please refer to your
        download for exact file name.}

    In operating systems what does not use the device name scheme shown above,
    such as macOS, the procedure is much of the same:
    First identify the device node which represents your USB device
    (for macOS, use \icode{diskutil list} in the Terminal application),
    then run the two commands as shown above using \textit{superuser privilege}.

    \subsection{Windows}

    Since there is no easy way to use dd or identify device node in a standard Windows installation
    (especially in versions prior to Windows 10), we recommend using the \textbf{Rufus} tool to achieve the same effect.

    Please refer to the Rufus site for usage details.

    \section{Booting the LiveKit Media}
    
    Before you boot the LiveKit media, several precautions and procedures should be followed. Here are a few common considerations for your reference:
    
    \subsection{What devices are supported?}

    LiveKit supports booting on the following device types:

    \begin{compactitem}
        \item BIOS- or (U)EFI-based x86\_64 and i486 devices.
        \item (U)EFI-based, or SBSA and SystemReady ES certified AArch64 devices.
        \item OpenFirmware- or Petitboot-based PowerPC Macintosh, POWER, and OpenPOWER devices.
        \item PMON- or Kunlun-based Loongson 2/3 Devices.
    \end{compactitem}

    \subsection{Do I have enough RAM?}
    
    For booting a Live media, assume that your RAM is at least half the size of the iso downloaded. Insufficient RAM will result in failure during start up or during system installation.

    While this is a less common issue with newer devices with RAM sizes in the gigabytes, devices targetted by AOSC OS/Retro, such as a Pentium MMX computer with 32MiB of RAM could easily run into this issue.

    \subsection{Other necessary tweaks for Secure Boot-enabled devices}

    LiveKit does not, and will probably not support Secure Boot.
    Therefore, for UEFI-based systems with Secure Boot support,
    you must disable Secure Boot in order to boot LiveKit.

    \section{Installing AOSC OS}
    
    After starting up LiveKit, you should be greeted with a command prompt with minimal direction shown on screen. To install AOSC OS, first connect to the Internet with NetworkManager. Simply run the nmtui command and you will be greeted with a simple configuration interface, follow the on-screen directions to configure your network connection.

    After you are connected, you may now use the included DeployKit installer to install AOSC OS. Simply run the following command to start the DeployKit installer:

\codeblock{
deploytkit
}

    You should be greeted with a wizard-like installer, simply follow the on-screen directions to install AOSC OS. This can take between minutes and nearly an hour depending on your device performance.
    
    
    \section{Maintaining or Repairing AOSC OS}

    In case of an unbootable system, or catastrophic system failures, LiveKit could prove useful as an environment for system maintenance and repair. Here are several classic use cases, for your reference.

    \subsection{Repairing a Broken ext4 System Partition}

    In case of a broken ext4 system partition, you may experience file access errors and system boot failure
    as the Kernel or initialisation mechanism fail to mount the system partition.
    Depending on your system configuration, you may not be able to repair your system partition
    without the help of an external environment.

    In this case, you may start LiveKit and, as you get to the command prompt,
    you may use the \icode{fsck.ext4} command to repair your ext4 partition.

    \subsubsection*{For Partition on IDE/ATA/SCSI/SAS Devices}
\codeblock{
fsck.ext4 -F /dev/sdXY
}

    \codenote{NOTE: Where X is the device number, denoted with a letter; Y is the partition number, denoted with a number.}

    \subsubsection*{For Partitions on NVMe-based Devices}
\codeblock{
fsck.ext4 -F /dev/nvmeXnYpZ
}

    \codenote{NOTE: Where X represents the device node, Y represents the device number, and Z represents the partition number, all denoted with numbers.}

    \subsubsection*{For Partitions on eMMC or SD Cards}
\codeblock{
fsck.ext4 -F /dev/mmcblkXpY
}

    \codenote{NOTE: Where X represents the device number, and Z represents the partition number, both denoted with numbers.}




    \subsection{Continuing an Interrupted System Update}

    In case of power loss, system crash, or other fatal errors during system updates,
    or user errors, such as powering down or interrupting system updates,
    you may run into boot or startup failures as a result.

    To attempt a repair, first identify your system partition by using the \icode{fdisk -l} command,
    and mount the system partition, the command differs based on your storage configuration.

    \subsubsection*{For Partition on IDE/ATA/SCSI/SAS Devices}
\codeblock{
mount /dev/sdXY /mnt
}

    \codenote{NOTE: Where X is the device number, denoted with a letter; Y is the partition number, denoted with a number.}

    \subsubsection*{For Partitions on NVMe-based Devices}
\codeblock{
mount /dev/nvmeXnYpZ /mnt
}

    \codenote{NOTE: Where X represents the device node, Y represents the device number, and Z represents the partition number, all denoted with numbers.}

    \subsubsection*{For Partitions on eMMC or SD Cards}
\codeblock{
mount /dev/mmcblkXpY /mnt
}

    \codenote{NOTE: Where X represents the device number, and Z represents the partition number, both denoted with numbers.}


    Next, you should connect to the Internet in case of need to download extra packages as a part of the repair procedures. Simply run the nmtui command and you will be greeted with a simple configuration interface, follow the on-screen directions to configure your network connection.

    After you have successfully connected to your specified network, you may now switch into your AOSC OS installation by running the following command:

\codeblock{
arch-chroot /mnt
}


    If you are able to see the AOSC OS command prompt,
    please run the following commands in succession,
    \textit{only if the current command completes successfully}.
    You should see a command prompt ending with \icode{\#} with a successfully completed command (otherwise,
    in case of an error, you will see \icode{!}; in case of missing commands or files, you will see \icode{?}):

\codeblock{
dpkg --configure -a\\
apt -f install\\
apt full-upgrade
}

    After the three commands completed successfully, use the following command to exit out of your AOSC OS installation:

\codeblock{
exit
}

    Further, run the following command to ensure all data are written to the disk:

\codeblock{
sync
}

    Finally, to reboot to your AOSC OS installation, execute the \icode{reboot} command.





    \subsection{Repairing Broken GRUB Installations}

    In case of an unbootable system as a result of a broken GRUB installation (this can happen as a result of an overwritten boot sector, accidentally deleted GRUB configuration, etc.), you will first need to mount your system partition:

    \subsubsection*{For Partition on IDE/ATA/SCSI/SAS Devices}
\codeblock{
mount /dev/sdXY /mnt
}

    \codenote{NOTE: Where X is the device number, denoted with a letter; Y is the partition number, denoted with a number.}

    \subsubsection*{For Partitions on NVMe-based Devices}
\codeblock{
mount /dev/nvmeXnYpZ /mnt
}

    \codenote{NOTE: Where X represents the device node, Y represents the device number, and Z represents the partition number, all denoted with numbers.}

    \subsubsection*{For Partitions on eMMC or SD Cards}
\codeblock{
mount /dev/mmcblkXpY /mnt
}

    \codenote{NOTE: Where X represents the device number, and Z represents the partition number, both denoted with numbers.}

    You may now switch into your AOSC OS installation by running the following command:

\codeblock{
arch-chroot /mnt
}

    If you are able to see the AOSC OS command prompt, we may now commence on repairing your GRUB installation.

    \subsubsection*{For (U)EFI- or Kunlun-based Systems}

    Further mount your EFI system partition. You may identify this partition,
    labelled as "EFI System," with the \icode{fdisk -l} command. After idenfying your EFI system partition,
    run the following command to mount it:

\codeblock{
mount \$DEVICE /efi
}

    \codenote{NOTE: \texttt{\$DEVICE} is the device node listed in the first column on the line corresponding
        the EFI system parititon (a parition of the "EFI System" type).}

    Execute the following commands in succession to repair your GRUB installation.
    \textit{Proceed only if the current command completes successfully}.
    You should see a command prompt ending with \icode{\#} with a successfully completed command
    (otherwise, in case of an error, you will see \icode{!}; in case of missing commands or files, you will see \icode{?}):

\codeblock{
grub-install --efi-directory=/efi \textbackslash\\
    --bootloader-id="AOSC OS"\\
grub-mkconfig -o /boot/grub/grub.cfg
}

    \subsubsection*{For BIOS-based Systems}

    First, identify the disk on which you installed AOSC OS with fdisk -l. This will be listed in a line starting with Disk.

    Execute the following commands in succession to repair your GRUB installation.
    \textit{Proceed only if the current command completes successfully}.
    You should see a command prompt ending with \icode{\#} with a successfully completed command
    (otherwise, in case of an error, you will see \icode{!}; in case of missing commands or files, you will see \icode{?}):

\codeblock{
grub-install \$DEVICE\\
grub-mkconfig -o /boot/grub/grub.cfg
}

    \codenote{NOTE: \texttt{\$DEVICE} is the device node listed in the first column on the line corresponding
        the EFI system parititon (a parition of the "EFI System" type).}

    \subsubsection*{For OpenFirmware-based Systems (Macintosh)}

    Execute the following commands in succession to repair your GRUB installation.
    \textit{Proceed only if the current command completes successfully}.
    You should see a command prompt ending with \icode{\#} with a successfully completed command
    (otherwise, in case of an error, you will see \icode{!}; in case of missing commands or files, you will see \icode{?}):

\codeblock{
grub-install\\
grub-mkconfig -o /boot/grub/grub.cfg
}


    \section{Further Help Needed?}

    If you would need further assistance with installing AOSC OS or repairing your AOSC OS installation,
    please contact us via Discord, Telegram, or \#aosc on the Libera Chat IRC.
    Our developers will be ready to assist and advise to help you navigate this sticky situation.

    \section{Reporting Bugs}

    If you have identified a bug while using LiveKit, or any tools pre-installed with LiveKit, please file a bug report at:

\codeblock{
https://github.com/AOSC-Dev/aosc-os-abbs
}





    \chapter{Installation (ARM)}

    This guide applies to ARMv4, ARMv6, and ARMv7.
    
    \section{Blah blah}

    // TODO





    \chapter{Afternotes}
    
    \section{Appendix I: Copyright}

    This manual is published with the \underline{CC BY 4.0} license.

    The source code of this manual is available at:

\codeblock{https://github.com/AOSC-Dev/optical-media-box}

    \vspace{1ex}This is the list of authors, sorted by Unicode codepoint order.

    \begin{tcolorbox}[colframe=black,colback=white]
        \begin{compactitem}
            \item Mingcong Bai
            \item Neruthes
        \end{compactitem}
    \end{tcolorbox}

    \section{Appendix II: Support}

    \begin{tabu} {lX}
        Website & \icode{\footnotesize https://aosc.io/} \\
        Telegram & \icode{\footnotesize https://t.me/joinchat/BMnG9079MOteeIJLxrW\_Mw} \\
        Discord & \icode{\footnotesize https://discord.gg/VYPHgt9} \\
        IRC & \icode{\footnotesize \#aosc} \\
    \end{tabu}

    % \begin{minipage}{0.35\textwidth}
    %     Hello World!
    % \end{minipage}\hfill%
    % \begin{minipage}{0.6\textwidth}
    %     Another column!
    % \end{minipage}




    % END OF CONTENT
    % \clearpage
    % \pagestyle{empty}

    % \makebackcoverpage
\end{document}